
% Planificacao da Aula 1

\begin{center}
{\large\bfseries\color{azulPDMI}Quantos triângulos se escondem num polígono?}\\[3pt]
{\small Planificação da aula - Matemática, 7.º ano}
\end{center}

\section*{1. Objeto}
\textbf{Tema:} Geometria. \textbf{Tópico:} Soma das amplitudes dos ângulos internos de um polígono convexo.\\
\textbf{Objeto matemático central:} dedução e justificação da soma das amplitudes dos ângulos internos de um polígono convexo com $n$ lados.\\
\textbf{Proposta de duração:} 50 minutos.

\section*{2. Objetivos}
\subsection*{Objetivos gerais}
A presente aula deve ajudar a desenvolver as capacidades de:
\begin{itemize}[leftmargin=*,itemsep=3pt]
\item \textbf{Raciocínio matemático.}
\item \textbf{Representações matemáticas:} usar a representação algébrica e a representação gráfica e fazer conexões entre estas representações.
\item \textbf{Trabalho colaborativo:} trabalhar em grupo.
\item \textbf{Comunicação matemática:} comunicar conteúdos matemáticos.
\end{itemize}

\subsection*{Objetivos específicos}
No final da aula, o aluno deve ser capaz de:
\begin{enumerate}[label=\textbf{OE\arabic*.},leftmargin=*,itemsep=4pt]
\item Decompor um polígono convexo em triângulos não sobrepostos, através do desenho de diagonais que não se intersetam.
\item Reconhecer a relação algébrica entre o número de lados ($n$) de um polígono convexo e o número de triângulos ($m$) em que se decompõe $\left[\,m=n-2\,\right]$.
\item Reconhecer a relação algébrica entre o número de lados ($n$) de um polígono convexo e a soma ($S$) das amplitudes dos seus ângulos internos $\left[\,S=(n-2)\times180^\circ\,\right]$.
\end{enumerate}

\section*{3. Pré-requisitos}
Ao iniciar a aula, o aluno deverá:
\begin{itemize}[leftmargin=*,itemsep=2pt]
\item Ter presente os conceitos de:
\begin{itemize}[itemsep=1pt]
\item vértices consecutivos ou adjacentes;
\item polígono convexo.
\end{itemize}
\item Ter presente que a soma das amplitudes dos ângulos internos de um triângulo é $180^\circ$.
\item Ser capaz de resolver equações simples.
\end{itemize}

\section*{4. Recursos e materiais}
\begin{itemize}[leftmargin=*,itemsep=2pt]
\item \textbf{A fornecer pelo professor:} folha em papel com o enunciado, a distribuir a cada grupo de trabalho.
\item \textbf{A trazer pelo aluno:}  lápis e borracha.
\item \textbf{Opcional:} régua e papel de rascunho.
\end{itemize}

\section*{5. Desenvolvimento da aula (50 minutos)}
\begingroup\small
\renewcommand{\arraystretch}{1.24}
\begin{longtable}{|p{1.35cm}|p{2.25cm}|p{5.55cm}|p{5.1cm}|}
\hline
\rowcolor{azulClaro}\textbf{Tempo}&\textbf{Fase / questões}&\textbf{Intervenção do professor}&\textbf{Atividade dos alunos}\\\hline
\endfirsthead
\hline\rowcolor{azulClaro}\textbf{Tempo}&\textbf{Fase / questões}&\textbf{Intervenção intencional do professor}&\textbf{Atividade dos alunos}\\\hline
\endhead
\textbf{0--5 min} (5 min)&Introdução. &
Apresenta a tarefa (tema e tópico). Recorda que a soma dos ângulos internos de um triângulo é 180º. Relembra o conceito de figura geométrica convexa. Coordena a organização a pares. &Leem o desafio, recordam que a soma dos ângulos internos de um triângulo é $180^\circ$ e interpretam «sem medir nenhum ângulo».\\\hline
\textbf{5--13 min} (8 min)&Parte I; 1.1&Circula e escuta antes de intervir. Regista as estratégias (tentativas, recorte de ângulos, decomposição em triângulos). Pergunta «O que já sabes que te possa ajudar?» sem sugerir a diagonal.&Procuram, a pares, uma estratégia para obter a soma dos ângulos internos do quadrilátero e do pentágono sem transferidor; desenham sobre as figuras e registam o raciocínio.\\\hline
\textbf{13--23 min} (10 min)&Parte I; 1.2, caixa de estratégia e 1.3&Introduz «diagonal» só depois da comparação. Pergunta «Há sobreposições ou regiões por preencher?». Apoia quem não relaciona os ângulos dos triângulos com os do polígono pedindo que os assinale com cores. A 1.3 fica essencialmente oral.&Comparam estratégias com o colega; traçam diagonais no quadrilátero e no pentágono; calculam $2\times180^\circ$ e $3\times180^\circ$; explicam, oralmente ou em registo curto, porque se podem somar os ângulos dos triângulos.\\\hline
\textbf{23--32 min} (9 min)&Parte II; 2.1--2.3&Insiste em traçar as diagonais a partir de um só vértice. Se surgirem diagonais cruzadas, pergunta «Os triângulos ficam sem sobreposições?». Pede: «O que representa $n-2$?» «Porquê $180^\circ$?». Não corrige de imediato o erro de $n$ triângulos.&Completam a tabela (quadrilátero e pentágono já estudados; hexágono novo), descrevem a relação entre lados e triângulos e propõem $S=(n-2)\times180^\circ$, explicando os fatores.\\\hline
\textbf{32--38 min} (6 min)&Parte II; 2.4&Pergunta «Quantos vértices não podes ligar ao vértice escolhido?» (ele próprio e os dois adjacentes). Distingue verificar na tabela de justificar para $n$ qualquer.&Justificam que a regra vale para qualquer polígono convexo, pensando nas $n-3$ diagonais a partir de um vértice e nos $n-2$ triângulos.\\\hline
\textbf{38--46 min} (8 min)&Discussão coletiva&Sequencia: estratégia por tentativa; decomposição do quadrilátero; tabela e regularidade; justificação com $n-3$ e $n-2$. Pede que expliquem cada fator antes de simplificar. «E se não fosse convexo?» fica como ponto de reflexão, sem aprofundar.&Apresentam 2 ou 3 resoluções previamente selecionadas; comparam estratégias; discutem o significado de $n-2$ e porque a decomposição funciona num polígono convexo.\\\hline
\textbf{46--50 min} (4 min)&Sistematização e conclusão&Institucionaliza a regra e a condição de convexidade. Recolhe as respostas individuais e identifica o que retomar na aula seguinte.&Completam individualmente «Para concluir» e determinam a soma para 12 lados ($1800^\circ$), justificando.\\\hline
\end{longtable}
\endgroup

\section*{6. Antecipação das dificuldades e respostas do professor}
\begingroup\small
\renewcommand{\arraystretch}{1.22}
\begin{longtable}{|p{2.55cm}|p{5.7cm}|p{6.65cm}|}
\hline
\rowcolor{azulClaro}\textbf{Questão / obstáculo}&\textbf{Produção ou dificuldade previsível}&\textbf{Intervenção docente planeada}\\\hline
\endfirsthead
\hline\rowcolor{azulClaro}\textbf{Questão / obstáculo}&\textbf{Produção ou dificuldade previsível}&\textbf{Intervenção docente planeada}\\\hline
\endhead
1.1: estratégia inicial&Medir os ângulos com o transferidor, apesar do pedido; recortar e juntar os ângulos; ou dizer que «não se consegue sem medir». Há ainda quem proponha uma regra sem fundamento (por exemplo, $360^\circ$ para qualquer figura).&«O que já sabes sobre a soma dos ângulos de um triângulo?» «Consegues ver triângulos dentro destas figuras?» Valorizar o recorte de ângulos como ideia, sem validar $360^\circ$ para todos. Não sugerir ainda a diagonal.\\\hline
1.1: desenho sobre as figuras&Traçar segmentos que não unem vértices, ou desenhar triângulos que não ocupam toda a figura, deixando regiões por preencher.&«Os extremos dos teus segmentos são vértices?» «Toda a região do polígono ficou coberta?» Pedir que marquem os triângulos com cores diferentes.\\\hline
1.2: comparação de estratégias&Os pares trocam resultados sem justificar; ou consideram que estratégias diferentes que dão o mesmo valor são equivalentes sem as relacionar.&«Porque é que as duas estratégias dão o mesmo resultado?» Pedir que cada elemento do par explique a ideia do colega. Registar as estratégias para a discussão.\\\hline
Caixa de estratégia, a)--b)&Traçar a diagonal do quadrilátero e contar 2 triângulos; no pentágono, traçar diagonais de vértices diferentes, que se cruzam, e obter mais de 3 triângulos; ou confundir diagonal com lado.&Recordar a definição: «Une vértices consecutivos?» Pedir o regresso a um único vértice e a contagem de triângulos sem sobreposições. Mostrar, se necessário, um exemplo com diagonais cruzadas e perguntar se os triângulos se sobrepõem.\\\hline
Caixa de estratégia, c)&Calcular $2\times180^\circ=360^\circ$ e $3\times180^\circ=540^\circ$ e não perceber que se trata da soma dos ângulos internos do polígono; ou somar $180^\circ$ por cada vértice (obter $720^\circ$ no quadrilátero).&«Que ângulos estão dentro de cada triângulo?» «Eles coincidem com os do quadrilátero?» Pedir que marquem, com a mesma cor, os ângulos do triângulo e a parte correspondente do ângulo do polígono.\\\hline
1.3: justificação da soma&Responder «porque dá certo» ou «porque o professor disse»; ou não reconhecer que, num vértice, o ângulo do polígono se divide em ângulos de dois triângulos.&«Há ângulos de triângulos que ficaram de fora ou que se repetem?» «O que acontece no vértice de onde partem as diagonais?» Aceitar respostas orais e registos curtos, e exigir referência às sobreposições e regiões por preencher.\\\hline
2.1: preenchimento da tabela&Traçar diagonais de vértices diferentes no hexágono; contar $n$ triângulos (6 no hexágono) em vez de $n-2$; esquecer-se de calcular a soma ou escrever $4\times180^\circ=620^\circ$ por erro de cálculo.&«De onde partem todas as diagonais?» «Quantos vértices não consegues ligar ao vértice escolhido?» Confrontar com o desenho: contar os triângulos um a um. Não indicar o valor; pedir que verifiquem a multiplicação.\\\hline
2.2: relação lados--triângulos&Dizer apenas «aumenta um de cada vez»; ou descrever a relação de forma recursiva («soma-se mais um triângulo») sem relacionar $n$ com o número de triângulos.&«Se tiver 10 lados, quantos triângulos terá?» Pedir que a descrição permita responder sem desenhar. Valorizar a ideia de «menos dois» e pedir que a justifiquem pelo desenho.\\\hline
2.3: simbolização&Escrever $S=n\times180^\circ$; escrever $(n-2)\times180^\circ$ por padrão, sem interpretar os fatores; ou escrever $180^\circ-2n$ ou $n-2\times180^\circ$ (sem parênteses).&«O que representa $n-2$?» «O que representa $180^\circ$?» Testar a expressão no quadrilátero e no pentágono. Se faltarem parênteses, pedir que calculem com $n=5$ e comparem com o valor da tabela.\\\hline
2.4: generalização&«Funciona no triângulo, quadrilátero, pentágono e hexágono, portanto funciona sempre»; ou explicar apenas com os exemplos da tabela, sem argumentar para $n$ qualquer.&«Quantos exemplos estudámos?» «Que argumento não depende do número de exemplos?» «Quantos vértices não podes ligar ao vértice escolhido?» (ele próprio e os dois adjacentes, logo $n-3$ diagonais e $n-2$ triângulos).\\\hline
2.4: convexidade&Não perceber porque o polígono tem de ser convexo; ou pensar que a regra só vale para polígonos regulares.&«Num polígono não convexo, as diagonais ficam sempre dentro do polígono?» Desenhar um exemplo não convexo como ponto de reflexão, sem aprofundar. Esclarecer que a regra não exige lados ou ângulos iguais.\\\hline
Conclusão e 12 lados&Copiar a expressão sem completar o «porque»; ou calcular $12\times180^\circ=2160^\circ$, esquecendo-se de subtrair 2; ou desenhar o dodecágono em vez de aplicar a regra.&«Quantos triângulos terias com 12 lados?» «Onde está esse número na tua expressão?» Pedir que expliquem o «porque» em linguagem própria. Recolher as respostas individuais para identificar o que retomar.\\\hline
Dificuldade técnica e de gestão&Material desenhado em escala pequena, dificultando o traçado; poucos alunos a terminar a tabela; discussões longas no trabalho a pares que reduzem o tempo da discussão coletiva.&Ter cópias adicionais das figuras e cores disponíveis; admitir que 1.3 e 2.4 sejam orais; avisar o tempo restante; preservar os últimos 4 minutos para o registo individual.\\\hline
\end{longtable}
\endgroup

\section*{7. Avaliação das aprendizagens}
\textbf{Instrumentos:} observação durante o trabalho a pares, respostas escritas no enunciado (incluindo «Para concluir»), argumentação na discussão coletiva e resposta individual ao caso dos 12 lados.

\textbf{Critério de resposta completa:} o aluno explicita que, a partir de um vértice, se traçam $n-3$ diagonais, que ficam $n-2$ triângulos, cada um com soma de ângulos internos $180^\circ$, pelo que $S=(n-2)\times180^\circ$; e aplica-o a 12 lados, obtendo $1800^\circ$. Respostas que apresentem apenas a expressão correta, sem justificação, revelam aprendizagem procedimental ainda incompleta.